\documentclass[10pt]{article}

% Specify the margins and text size.
\setlength{\textwidth}{6.6in}
\setlength{\textheight}{9.5in}
\setlength{\oddsidemargin}{0pt}
\setlength{\evensidemargin}{0pt}
\setlength{\topmargin}{0pt}
\setlength{\hoffset}{.05in}
\setlength{\voffset}{-.8in}

\setlength{\parskip}{5pt}
\setlength{\parindent}{0pt}

% Load some fonts and symbol packages
\usepackage{latexsym}
\usepackage{pifont}       % contains 'star' symbol for counterinsurgency handbook title
\usepackage{yfonts} 
\usepackage{amsmath}
\usepackage{amsfonts}

\usepackage{graphicx}     % actually, this is loaded by pstricks
\usepackage[T1]{fontenc}
\usepackage{ifthen}
\usepackage{pstricks,pst-func,pst-text,pst-node,multido,pst-plot,calc,pst-3dplot}
%\usepackage[all]{xy}
%\usepackage{animate}

% The hyperref package inserts links into the pdf.
\definecolor{MyLinkColor}{rgb}{.1,.2,1}
\definecolor{MyCiteColor}{rgb}{.1,1,.2}
\definecolor{MyURLColor}{rgb}{.4,.4,.4}
\usepackage[backref=true,pagebackref=false,hyperindex,colorlinks=true,
  linkcolor=MyLinkColor,urlcolor=MyURLColor]{hyperref}

% The tweaklist package is something I found on the web.  It provides a simple interface
% for making changes to spacing used in the itemize and enumerate environments.  Comment
% this out if you don't care to use tweaklist.
\usepackage{tweaklist}
\renewcommand{\itemhook}{\setlength{\parskip}{2pt}\setlength{\parsep}%
{1pt}\setlength{\topsep}{0pt}\setlength{\itemsep}{0pt}}

\newcommand{\N}{\ensuremath\mathbb{N}}
\newcommand{\Q}{\ensuremath\mathbb{Q}}
\newcommand{\R}{\ensuremath\mathbb{R}}
\newcommand{\C}{\ensuremath\mathbb{C}}
\newcommand{\z}{\hphantom{0}}

\begin{document}
\pagestyle{empty}

\href{http://www.su.edu}{\includegraphics[height=1.75cm]{sulogo.eps}}
\vspace{-1.69cm}

{\small{\ }\hfill
\begin{tabular}{cl}
& Math 207\\& Introduction to Statistics\\
\end{tabular}
}
\setlength{\baselineskip}{1.05\baselineskip}
\medskip

{\setlength{\baselineskip}{1.05\baselineskip}
\begin{center}
\textbf{\large  Worksheet:  Correlation}
\end{center}

\newcommand{\Z}{\hphantom{0}}
1. Two lists of data $x_1$, $\dots$, $x_n$ and $y_1$, $\dots$, $y_n$ have
the following properties.  
\begin{center}
\begin{tabular}{ccc}
$\mbox{mean}_x = 10$ & $\mbox{sd}_x=2$ & \\
$\mbox{mean}_y = 3\Z$  & $\mbox{sd}_y=4$ & $r=0.5$
\end{tabular}
\end{center}

\hspace{20pt} a) Find the equation for the regression line and make a sketch.
\vspace{1.75in}

\hspace{20pt} b) Predict the $y$ value if $x=8$.  If $x=14$?
\vfill


\href{http://www.su.edu}{\includegraphics[height=1.75cm]{sulogo.eps}}
\vspace{-1.69cm}

{\small{\ }\hfill
\begin{tabular}{cl}
& Math 207\\ & Introduction to  Statistics\\
%                & Worksheet\\
\end{tabular}
}
\begin{center}
\textbf{\large  Worksheet:  Correlation}
\end{center}

1. Two lists of data $x_1$, $\dots$, $x_n$ and $y_1$, $\dots$, $y_n$ have
the following properties.  
\begin{center}
\begin{tabular}{ccc}
$\mbox{mean}_x = 10$ & $\mbox{sd}_x=2$ & \\
$\mbox{mean}_y = 3\Z$  & $\mbox{sd}_y=4$ & $r=0.5$
\end{tabular}
\end{center}

\hspace{20pt} a) Find the equation for the regression line and make a sketch.
\vspace{1.75in}

\hspace{20pt} b) Predict the $y$ value if $x=8$.  If $x=14$?
\vspace{.2in}

\eject
\newcommand{\Zc}{\hphantom{$\;$}}
2. For the given data, calculate 
$\mbox{mean}_x$, $\mbox{sd}_x$,
$\mbox{mean}_y$, $\mbox{sd}_y$ and  $r$ then find the equation for
the regression line.

\begin{tabular}{cr}
$x$ & $y$\\\hline
1   & 3\\
2   & 1\\
6   &   -1\\
\end{tabular}
\vfill
2. For the given data, calculate 
$\mbox{mean}_x$, $\mbox{sd}_x$,
$\mbox{mean}_y$, $\mbox{sd}_y$ and  $r$ then find the equation for
the regression line.

\begin{tabular}{cr}
$x$ & $y$\\\hline
1   & 3\\
2   & 1\\
6   &   -1\\
\end{tabular}
\vspace{3in}
\eject

\end{document}
